% META: {"id": "TSPE-GEOESPACE-ME-006", "chapitre": "TSPE-GEOMETRIE-ESPACE", "type_objet": "methode", "capacites_codes": ["C8", "C11", "C14"], "methodes": ["M6"], "sources_inspiration": [], "status": "generated"}
% BEGIN-VERIFY
% from sympy import Matrix, symbols, solve, Eq, sqrt, simplify
% t = symbols('t')
% n=Matrix([1,1,1]); d=-3; M=Matrix([3,3,3])
% H = M + solve(Eq(n.dot(M + t*n) + d, 0), t)[0]*n
% assert list(H) == [1,1,1]
% assert n.dot(H) + d == 0
% assert simplify((H-M).norm() - 2*sqrt(3)) == 0
% assert simplify((H-M).norm() - abs(n.dot(M)+d)/n.norm()) == 0
% END-VERIFY
\begin{fichemethode}{M6}{Projeter orthogonalement : coordonnées du projeté et distance}

\quandUtiliser{%
  Pour déterminer les coordonnées du projeté orthogonal d'un point sur un plan,
  pour en déduire la distance de ce point au plan, et pour justifier que ce
  projeté est le point du plan le plus proche. Les trois questions suivent la
  même procédure.
}

\pasApas{%
  \begin{enumerate}
    \item \textbf{Lire} le vecteur normal $\vec{n}(a;b;c)$ dans l'équation
          $ax+by+cz+d=0$ du plan.
    \item \textbf{Paramétrer} la droite passant par $M$ et dirigée par
          $\vec{n}$~: $X=M+t\,\vec{n}$.
    \item \textbf{Reporter} ces coordonnées dans l'équation du plan et
          résoudre en $t$~: la solution est unique.
    \item \textbf{Calculer} $H$ en reportant cette valeur de $t$.
    \item \textbf{En déduire} la distance $MH=\lVert\vec{MH}\rVert$, qui vaut
          aussi
          $\dfrac{\lvert a x_M+b y_M+c z_M+d\rvert}{\sqrt{a^2+b^2+c^2}}$.
    \item \textbf{Pour justifier} que $H$ est le plus proche~: pour tout $N$
          du plan distinct de $H$, $\vec{MH}\cdot\vec{HN}=0$, donc
          $MN^2=MH^2+HN^2>MH^2$ par Pythagore.
  \end{enumerate}
}

\verifier{%
  Contrôler que $H$ vérifie bien l'équation du plan. Contrôler ensuite que les
  deux calculs de la distance --- par la longueur $MH$ et par la formule ---
  donnent le même résultat~: ils se contrôlent mutuellement.
}

\pieges{%
  Annuler une coordonnée ne projette que sur un plan de coordonnées~; pour un
  plan quelconque il faut suivre la normale. Dans la formule de la distance,
  ne pas oublier la valeur absolue au numérateur ni la norme au dénominateur.
  Enfin, l'inégalité $MN>MH$ est \emph{stricte} et suppose $N\neq H$.
}

\exempleRedige{%
  $\mathcal{P}:\ x+y+z-3=0$ et $M(3;3;3)$. Ici $\vec{n}(1;1;1)$, d'où
  $X=(3+t;3+t;3+t)$. En reportant~: $3(3+t)-3=0$, soit $3t+6=0$ et $t=-2$.
  Donc $H(1;1;1)$, qui vérifie $1+1+1-3=0$. Enfin
  $MH=\sqrt{(-2)^2\times3}=2\sqrt3$, et la formule donne
  $\dfrac{\lvert 9-3\rvert}{\sqrt3}=\dfrac{6}{\sqrt3}=2\sqrt3$~: les deux
  concordent.
}

\end{fichemethode}
